

%----------------------------------------------------------------------------------------
%	NAME AND CONTACT INFORMATION SECTION
%----------------------------------------------------------------------------------------


% Personal Data
\firstname{David}
\familyname{Molik}
\title{Evaluation - 2025}
\address{University Library}{Kansas State University}
%\mobile{+1 (716) 982-5235, +1 (785) 532-5399}
%\email{molik@ksu.edu} 
%\homepage{http://david.molik.co}{david.molik.co}
%\extrainfo{
  %\faGithub\href{https://github.com/molikd}{ molikd} \quad
  %\faLinkedin\href{https://www.linkedin.com/in/dmolik/}{ dmolik} \quad
  %\aiOrcid\href{https://orcid.org/0000-0003-3192-6538}{ 0000-0003-3192-6538}
%}

%\social[linkedin][www.linkedin.com]{name}
% The first argument is %the url for the clickable link, the second argument is the url displayed in the %template - this allows special characters to be displayed such as the tilde in this %example

%\photo[70pt][0.3pt]{picture} % The first bracket is the picture height, the second is %the thickness of the frame around the picture (0pt for no frame)
%\quote{Not Attention, Patience is all we need.}

%----------------------------------------------------------------------------------------


\portfoliotitle % Repeatable section header (see main.tex)

%\title{Personal Goals (2025--2030)}

%----------------------------------------------------------------------------------------
%   EVALUATION
%----------------------------------------------------------------------------------------


\section{Introduction}

I started at Kansas State Libraries in February of this year. Looking back on this evaluation period, I would characterize my work as strategic, intentionally ambitious, and necessarily iterative. The Center for Scholarly Publishing---now the Open Publishing Exchange (OPE)---represents a broad and evolving portfolio spanning scholarly publishing, emerging data services, professional development, and institutional strategy. Ideally such an organization should be able to affect real change in both Kansas State's research enterprise, as well as be a force in Student Success. This year marked a period of transition for both the department and the Libraries, and while this breadth occasionally required balancing depth against momentum, the overall impact has been positive in positioning OPE for its next phase of growth.

\vspace{\baselineskip}

During the evaluation period, I met and exceeded expectations outlined in my Unclassified Professional Plan (UPP) through sustained leadership in directed service, meaningful engagement in non-directed service, and continued scholarly and creative activity aligned with my role as Department Head. In parallel with administrative responsibilities, I maintained active involvement in scholarly publishing, research dissemination, and standards development that directly support the mission of the Libraries and the University.

\vspace{\baselineskip}

One of the most significant lessons I learned this year is the importance of pacing and structural clarity when leading change. Initiatives such as launching a new journal, reshaping open textbook governance, and advancing data curation capacity all benefited from early vision-setting, but also revealed how essential it is to clearly articulate decision points, ownership, and sustainability early in the process. I also gained a deeper appreciation for how much institutional change depends on informal relationship-building alongside formal structures.

\vspace{\baselineskip}

If I were to approach this year again, I would be more deliberate about narrowing the number of simultaneous new initiatives and more explicit about delegation at earlier stages; however, I also recognize that some breadth was inevitable given the department's scope and the Libraries' evolving needs. While I am comfortable working across domains, I recognize that long-term success of the department depends on ensuring that projects are not overly dependent on individual effort, and that faculty and staff expertise is surfaced and empowered as early as possible.

\vspace{\baselineskip}

Looking ahead, I would like to see several of these efforts transition from launch mode to shared ownership and refinement. In the future, I would prioritize bringing additional stakeholders into planning conversations earlier, particularly colleagues in Research IT, the Provost’s Office, and faculty partners, so that initiatives mature with broader institutional buy-in. I recognize that doing this while retaining the speed at which the department is able to innovate may be difficult, but my goal is to continue evolving from hands-on implementation toward a more facilitative leadership role that emphasizes alignment, sustainability, and capacity-building.

\vspace{\baselineskip}

As Department Head of the Open Publishing Exchange, I led a team of three Scholarly Communication librarians and authored a departmental landscape document that served as the basis for the department’s strategic direction. I established priorities focused on reorganizing the Open and Alternative Textbook Initiative, building a Research Data Management and Curation program, and leveraging library assets such as New Prairie Press to pursue university-level responses to the serials crisis. In addition, I provided information and analysis to Libraries leadership through a Virginia Tech Libraries site visit and report, attendance at the Center for Networked Information meeting, and liaison work with Research IT and the Provost’s Office related to Wildcat Scholar and Kansas State University’s CRIS/RIMS planning.

\vspace{\baselineskip}

As part of this evaluation cycle, I co-founded and launched \textit{Cyberbiosecurity Quarterly} as Co-Editor-in-Chief, establishing a new peer-reviewed scholarly venue hosted on New Prairie Press. In its inaugural year, I led a seven-person editorial team, contributed to editorial governance and workflow development, and personally reviewed multiple manuscript submissions. I also contributed to peer-reviewed research outputs during the evaluation period, including middle authorship on articles published in \textit{microPublication Biology} and \textit{Genetics} (``Guidelines for Gene and Genome Assembly Nomenclature''). While the majority of the \textit{Genetics} manuscript was completed prior to my current position, I finalized contributions and supported the work through publication during this evaluation cycle. In addition, I co-authored the \textit{System Security Assessment Plan Template}, a Research Data Alliance working group output released in November 2025. These scholarly and creative activities provide context for the service contributions described below and reflect sustained engagement with communities relevant to my UPP responsibilities.


%----------------------------------------------------------------------------------------
%   PROFESSIONAL CONDUCT
%----------------------------------------------------------------------------------------

\section{Professional Conduct}

Throughout the evaluation period, I interacted cooperatively with colleagues and campus partners, demonstrating civility, respect, and professionalism consistent with University Handbook C46.1. I showed respect for differing viewpoints in the exchange of ideas and followed appropriate supervisory guidance while exercising independent judgment in the execution of my responsibilities. I maintained a professional and collegial demeanor in all internal and external interactions.

%----------------------------------------------------------------------------------------
%   DIRECTED SERVICE
%----------------------------------------------------------------------------------------

\section{Directed Service}

\vspace{\baselineskip}

\subsection*{Core Responsibilities (High)}

During the evaluation year, I fulfilled and exceeded my core responsibilities as Department Head for the Open Publishing Exchange. I provided leadership and oversight for scholarly publishing, open educational resources, and research-support services, ensuring alignment with Libraries and University priorities. I balanced administrative, supervisory, and strategic responsibilities while representing the Libraries in campus-wide discussions related to scholarly communication, research infrastructure, and compliance.

\vspace{\baselineskip}

\subsection*{Complete K-State Supervisory Foundations (Med)}

I completed the required K-State Supervisory Foundations training, strengthening my capacity in personnel management, communication, and application of University policies. Knowledge gained through this training informed my approach to staff supervision, mentoring, and performance evaluation.

\vspace{\baselineskip}

In addition, I engaged in professional development and leadership training opportunities to strengthen my effectiveness as a Department Head. These included participation in Harvard’s Leadership Institute for Academic Librarians (LIAL),  the K-State Fall New Department Head Institute, and the University of Kansas Assessment in Action Workshop. Collectively, these experiences enhanced my capacity in strategic leadership, organizational assessment, personnel management, and evidence-based decision-making.

\vspace{\baselineskip}

A key lesson from these experiences was the importance of explicitly separating strategic direction from operational execution in supervisory contexts (i.e., when to think through a problem and when to act). As a result of these trainings and experiences, I became more deliberate about clarifying expectations, documenting decision rationales, and distinguishing between strategic conversations and implementation/working meetings. I applied these lessons by adjusting how I structure project check-ins and how I contextualize department- and project-level meetings.

\vspace{\baselineskip}


\subsection*{Build Data Curation and Management Program (High)}

I made substantial progress toward building a data curation and management program aligned with FAIR principles and federal open-access mandates. I served as a library liaison to the Office of Research on issues related to data repositories, compliance with the Nelson Memo, and CRIS/RIMS planning. This work positioned the Libraries as an active partner in campus research infrastructure planning.

\vspace{\baselineskip}

As part of this work, I conceived and organized the ``Tea with CSP'' discussion series, facilitating informal but substantive conversations on scholarly communication, open publishing, and research infrastructure. These discussions included participation by national leaders such as James Butcher and Natalie Meyers and supported professional development within the department.

\vspace{\baselineskip}

A key lesson from this work was that successful data curation and management programs depend as much on relationship-building and shared vocabulary as on technical infrastructure. Our early conversations revealed differing assumptions across units about ownership, scope, and responsibility for data services. In response, I have started to emphasize framing discussions around researcher workflows and compliance needs (i.e., thinking about the problem from a research lifecycle perspective, and not singularly data curation or single projects). 

\vspace{\baselineskip}

\subsection*{OPE Processes and Procedures Updating (Med)}

I led internal efforts to review and update CSP/OPE processes and procedures, including scholarly publishing workflows, editorial governance, and service models. This included alignment with emerging publishing platforms such as Fulcrum and Janeway and refinement of internal documentation to support consistent and sustainable operations.

\vspace{\baselineskip}

A key lesson from this work was the importance of making implicit workflows explicit. As processes are reviewed and updated, it has become clear that undocumented practices can create unnecessary friction and limit scalability. This reinforced the value of clear documentation and shared reference points to support consistency, continuity, and onboarding as the department’s scope continues to evolve.

\vspace{\baselineskip}

\subsection*{Build OER Program (High)}

I expanded and supported the Open and Alternative Textbook Initiative (OATI), contributing to increased faculty engagement in open and alternative instructional materials. A significant component of this work involved designing a new support structure for OER services, including drafting bylaws for a reorganized OATI committee. The revised committee structure will take effect in early 2026.

\vspace{\baselineskip}

A key lesson from this work was that sustained growth in OER programs depends on clear governance and shared ownership rather than ad hoc support. While faculty interest is critical, long-term success required formalizing decision-making structures and roles to ensure continuity, transparency, and institutional commitment. This insight has directly informed the design of the revised OATI committee structure.


%----------------------------------------------------------------------------------------
%   NON-DIRECTED SERVICE
%----------------------------------------------------------------------------------------

\section{Non-Directed Service}

\subsection*{Data Curation and Data Analysis Training (Med)}

I contributed to data curation and analysis training beyond my directed responsibilities by serving as an instructor for the Incoming Life Sciences Student Carpentries Program. I also participated in monthly Great Plains Network Carpentries Instructor meetings, contributing to curriculum discussion and regional coordination.

\vspace{\baselineskip}

A key lesson from this work was that sustainable data training depends on building shared instructional capacity rather than relying on individual expertise. Participation in regional instructor coordination highlighted the value of common curricula, peer feedback, and distributed teaching models. These experiences informed how I think about scaling data literacy efforts beyond single workshops or programs.


\vspace{\baselineskip}

\subsection*{EHS Committee (Low)}

I fulfilled my responsibilities as a member of the EHS Committee, participating as required and contributing to committee activities consistent with its scope and charge. As my service on this committee began during this evaluation period, this year focused primarily on orientation and gaining familiarity with the committee’s scope, processes, and priorities, which will inform deeper engagement in future cycles.


\vspace{\baselineskip}

\subsection*{Wildcat Scholar Committee (Med)}

I served on the Wildcat Scholar Committee, contributing scholarly communication and research infrastructure expertise to discussions related to research visibility, repository integration, and institutional research profiling. This year focused primarily on orientation and gaining familiarity with the committee’s scope, processes, and priorities, which will inform deeper engagement in future cycles.


%----------------------------------------------------------------------------------------
%   RESEARCH AND CREATIVE CONTEXT
%----------------------------------------------------------------------------------------

\section{Research and Creative Context}

In parallel with the service activities described above, I maintained engagement in scholarly and creative activities that support my UPP responsibilities. These activities complement and reinforce my directed and non-directed service contributions in scholarly communication, publishing infrastructure, and research data governance.

\vspace{\baselineskip}

A key lesson from maintaining this engagement was the value of remaining an active participant in scholarly communities while serving in an administrative role. Continued involvement in research and standards development helped ground strategic and operational decisions in current disciplinary practice, ensuring that service and infrastructure initiatives remained responsive to researcher needs.

\vspace{\baselineskip}

This reciprocal relationship between scholarship and service has been a defining feature of my approach during this evaluation period.
